\chapter{Структуры}

\section{FILE и структура \_IO\_FILE}

\begin{mylisting}{FILE /libio/bits/types/FILE.h}{code:file}{C}
#ifndef __FILE_defined
#define __FILE_defined 1
struct _IO_FILE;
/* The opaque type of streams.  This is the definition used elsewhere.  */
typedef struct _IO_FILE FILE;
#endif
\end{mylisting}

\begin{mylisting}{\_IO\_FILE /libio/bits/types/struct\_FILE.h}{code:iofile}{C}
struct _IO_FILE
{
	int _flags;		/* High-order word is _IO_MAGIC; rest is flags. */
	
	/* The following pointers correspond to the C++ streambuf protocol. */
	char *_IO_read_ptr;	/* Current read pointer */
	char *_IO_read_end;	/* End of get area. */
	char *_IO_read_base;	/* Start of putback+get area. */
	char *_IO_write_base;	/* Start of put area. */
	char *_IO_write_ptr;	/* Current put pointer. */
	char *_IO_write_end;	/* End of put area. */
	char *_IO_buf_base;	/* Start of reserve area. */
	char *_IO_buf_end;	/* End of reserve area. */
	
	/* The following fields are used to support backing up and undo. */
	char *_IO_save_base; /* Pointer to start of non-current get area. */
	char *_IO_backup_base;  /* Pointer to first valid character of backup area */
	char *_IO_save_end; /* Pointer to end of non-current get area. */
	
	struct _IO_marker *_markers;
	
	struct _IO_FILE *_chain;
	
	int _fileno;
	int _flags2:24;
	/* Fallback buffer to use when malloc fails to allocate one.  */
	char _short_backupbuf[1];
	__off_t _old_offset; /* This used to be _offset but it's too small.  */
	
	/* 1+column number of pbase(); 0 is unknown. */
	unsigned short _cur_column;
	signed char _vtable_offset;
	char _shortbuf[1];
	
	_IO_lock_t *_lock;
	#ifdef _IO_USE_OLD_IO_FILE
};

struct _IO_FILE_complete
{
	struct _IO_FILE _file;
	#endif
	__off64_t _offset;
	/* Wide character stream stuff.  */
	struct _IO_codecvt *_codecvt;
	struct _IO_wide_data *_wide_data;
	struct _IO_FILE *_freeres_list;
	void *_freeres_buf;
	struct _IO_FILE **_prevchain;
	int _mode;
	#if __WORDSIZE == 64
	int _unused3;
	#endif
	__uint64_t _total_written;
	#if __WORDSIZE == 32
	int _unused3;
	#endif
	/* Make sure we don't get into trouble again.  */
	char _unused2[12 * sizeof (int) - 5 * sizeof (void *)];
};
\end{mylisting}

\section{Структура stat}

%\vspace*{-2cm}
\begin{mylisting}{struct stat}{code:stat}{C}
struct stat {
	dev_t     st_dev;     /* устройство */
	ino_t     st_ino;     /* inode */
	mode_t    st_mode;    /* режим доступа */
	nlink_t   st_nlink;   /* количество жестких ссылок */
	uid_t     st_uid;     /* идентификатор пользователя-владельца */
	gid_t     st_gid;     /* идентификатор группы-владельца */
	dev_t     st_rdev;    /* тип устройства (если это устройство)*/
	off_t     st_size;    /* общий размер в байтах */
	blksize_t st_blksize; /* размер блока ввода-вывода в файловой системе */
	blkcnt_t  st_blocks;  /* количество выделенных 512-байтных блоков */
	union { time_t st_atime; struct timespec st_atim; }; /* время последнего доступа */
	union { time_t st_mtime; struct timespec st_mtim; }; /* время последней модификации */
	union { time_t st_ctime; struct timespec st_ctim; }; /* время последнего изменения */
};
\end{mylisting}

\newpage
\chapter{Программа №1}

\section{Код программы}

\begin{mylisting}{Код первой программы}{code:p1}{C}
#include <stdio.h>
#include <fcntl.h>

int main()
{
	int fd = open("alphabet.txt",O_RDONLY);
	
	FILE *fs1 = fdopen(fd,"r");
	char buff1[20];
	setvbuf(fs1,buff1,_IOFBF,20);
	
	FILE *fs2 = fdopen(fd,"r");
	char buff2[20];
	setvbuf(fs2,buff2,_IOFBF,20);
	
	int flag1 = 1, flag2 = 2;
	while(flag1 == 1 || flag2 == 1)
	{
		char c;
		flag1 = fscanf(fs1,"%c",&c);
		if (flag1 == 1) 
		{
			fprintf(stdout,"%c",c);
		}
		flag2 = fscanf(fs2,"%c",&c);
		if (flag2 == 1)
		{ 
			fprintf(stdout,"%c",c); 
		}
	}
	return 0;
}
\end{mylisting}

\section{Результат работы программы}

\img{p1}{Результат работы первой программы}{pic:p1}

\section{Анализ работы программы}

Системный вызов open открывает файл alphabet.txt в режиме <<только для чтения>>, создаёт структуру file. Ссылка на эту структуру помещается в таблицу файловых дескрипторов процесса. Поскольку дескрипторы 0, 1 и 2 заняты стандартными потоками ввода-вывода, системный вызов возвращает наименьший свободный дескриптор (3), который присваивается переменной fd.

Библиотечной функции fdopen передается дескриптор fd и режим доступа <<только для чтения>>. Функция возвращает указатель на инициализированную структуру FILE, где полю \_fileno присваивается значение fd.

Функция setvbuf задаёт размер буфера 20 байт для fs1 и fs2.

Функция fscanf для fs1 сначала считывает 20 символов (<<a>>---<<t>>) из файла "alphabet.txt" в буфер, при этом значение поля f\_pos структуры file увеличивается на 20. Затем из буфера извлекается первый символ (<<a>>) и записывается в переменную c, а указатель \_IO\_read\_ptr в структуре fs1 сдвигается на 1 байт. Вызов fprintf выводит этот символ в стандартный поток вывода.

Так как fs1 и fs2 ссылаются на один и тот же открытый файл, где значение f\_pos равняется 20, то при вызове fscanf для fs2 в буфер считываются оставшиеся 6 символов (<<u>>---<<z>>) из файла "alphabet.txt". Первый символ из буфера fs2 (<<u>>) записывается в переменную c и выводится в стандартный поток вывода.

В последующих итерациях цикла поочерёдно будут выводиться символы из буферов fs1 и fs2.
Когда будет достигнут конец буфера fs2 (сравняются указатели \_IO\_read\_ptr и \_IO\_read\_end), то программа продолжит выводить оставшиеся символы только из буфера fs1 (вплоть до t), пока не будет достигнут конец буфера fs1.

\section{Схема связи структур}

\img{1.pdf}{Схема связи структур в первой программе}{pic:scheme_1}

\newpage
\chapter{Программа №2}

\section{Однопоточный вариант программы}

\subsection{Код программы}

\begin{mylisting}{Код однопоточной первой программы}{code:p2_1thread}{C}
#include <fcntl.h>
#include <unistd.h>

int main()
{
	char c;
	int fl1 = 1, fl2 = 1; 
	
	int fd1 = open("alphabet.txt",O_RDONLY);
	int fd2 = open("alphabet.txt",O_RDONLY);
	
	while((fl1 == 1) && (fl2 == 1))
	{
		if ((fl1 = read(fd1,&c,1)) == 1)
		write(1,&c,1);
		if (fl1 == 1 && ((fl2 = read(fd2,&c,1)) == 1))
		write(1,&c,1);
	}
	return 0;
}
\end{mylisting}

\subsection{Результат работы программы}

\img{p2_read_1thread}{Результат работы однопоточной второй программы}{pic:p1}

\subsection{Анализ работы программы}

Системный вызов open создаёт два дескриптора открытого файла в режиме <<только для чтения>> (fd1, fd2). Соответственные дескрипторам структуры file ссылаются на один и тот же inode файла <<alphabet.txt>>. Ссылки на эти структуры помещаются в таблицу открытых процессом файлов под индексами 3 и 4, так как дескрипторы 0,1 и 2 заняты стандартными потоками.

Системный вызов read для дескриптора fd1 считывает один символ (<<a>>), что приводит к увеличению на 1 значения поля f\_pos в соответствующей структуре file. Считанный символ выводится в стандартный поток вывода.

Так как дескрипторам fd1 и fd2 соответствуют независимые структуры file, то для fd2 значение поля f\_pos равно 0. Поэтому при вызове read для fd2 будет считан тот же самый первый символ (<<a>>), который также выведется в стандартный поток вывода.

Поочерёдное чтение и вывод символов будут продолжаться в цикле, пока значения f\_pos в обеих структурах file не достигнут конца файла (размера файла, хранящегося в поле i\_size и равного 26 байтам).

Таким образом, в стандартный поток вывода каждый символ алфавита будет выведен строго последовательно дважды.

\section{Вариант программы с 1 дополнительным потоком}

\subsection{Код программы}

\begin{mylisting}{Код программы с 1 дополнительным потоком}{code:p2_2thread_join}{C}
#include <fcntl.h>
#include <pthread.h>
#include <unistd.h>
#include <stdio.h>
#include <stdlib.h>

void *thread_func(void *arg)
{
	char c;
	int fl = 1;
	int fd = *((int *)arg);
	
	while (fl == 1)
	{
		if ((fl = read(fd, &c, 1)))
		write(1, &c, 1);
	}
	return NULL;
}

int main()
{
	pthread_t t;
	char c;
	int fl = 1;
	
	int fd1 = open("alphabet.txt",O_RDONLY);
	int fd2 = open("alphabet.txt",O_RDONLY);
	
	if (pthread_create(&t, NULL, thread_func, &fd1) != 0)
	{
		perror("pthread_create");
		exit(1);
	}
	
	while (fl == 1)
	{
		if ((fl = read(fd2,&c,1)) == 1)
		write(1,&c,1);
	}
	pthread_join(t, NULL);
	return 0;
}
\end{mylisting}

\subsection{Результат работы программы}

\img{p2_read_2thread_join}{Результат работы программы с 1 дополнительным потоком}{pic:p1}

\subsection{Анализ работы программы}

Вызов pthread\_create создаёт дополнительный поток, которому в качестве аргумента передаётся указатель на дескриптор fd1. Поскольку на создание потока требуется время, главный поток может успеть прочитать и вывести либо весь алфавит, либо его часть.

После создания дополнительного потока чтение и вывод символов осуществляются обоими потоками в зависящей от планировщика последовательности.

Когда главный поток завершит чтение по своему дескриптору, он выйдет из цикла и заблокируется в вызове pthread\_join, ожидая завершения дополнительного потока. Дополнительный поток продолжит работу и выведет оставшуюся часть своих символов.

Таким образом, каждый символ файла будет напечатан дважды, однако точный порядок их вывода не определён и полностью зависит от того, как планировщик распределяет процессорное время между потоками.

\section{Вариант программы с 2 дополнительными потоками}

\subsection{Код программы}

\begin{mylisting}{Код программы с 2 дополнительными потоками}{code:p2_3thread_join}{C}
#include <fcntl.h>
#include <pthread.h>
#include <unistd.h>
#include <stdio.h>
#include <stdlib.h>

void *thread_func(void *arg)
{
	char c;
	int fl = 1;
	int fd = *((int *)arg);
	
	while (fl == 1)
	{
		if ((fl = read(fd, &c, 1)))
		write(1, &c, 1);
	}
	return NULL;
}

int main()
{
	pthread_t t1, t2;
	
	int fd1 = open("alphabet.txt",O_RDONLY);
	int fd2 = open("alphabet.txt",O_RDONLY);
	
	if (pthread_create(&t1, NULL, thread_func, &fd1) != 0)
	{
		perror("pthread_create");
		exit(1);
	}
	if (pthread_create(&t2, NULL, thread_func, &fd2) != 0)
	{
		perror("pthread_create");
		exit(1);
	}
	
	pthread_join(t1, NULL);
	pthread_join(t2, NULL);
	
	return 0;
}
\end{mylisting}

\subsection{Результат работы программы}

\img{p2_read_3thread_join}{Результат работы программы с 2 дополнительными потоками}{pic:p1}

\subsection{Анализ работы программы}

Два вызова pthread\_create создают два дополнительных потока. Каждому из них передаётся указатель на свой файловый дескриптор (первому --- fd1, второму --- fd2). Главный поток, создав рабочие потоки, сразу же блокируется в вызовах pthread\_join, ожидая их завершения.

На создание потоков требуется время и их порядок выполнения зависит от планировщика, поэтому порядок вывода символов неизвестен.

Таким образом, каждый символ файла будет напечатан дважды, однако точный порядок их вывода не определён и полностью зависит от того, как планировщик распределяет процессорное время между дополнительными потоками.

\section{Схема связи структур}

\img{2.pdf}{Схема связи структур во второй программе}{scheme_2}

\newpage
\chapter{Программа №3}

\section{Вариант программы с open, write и close}

\subsection{Код программы}

\begin{mylisting}{Код программы с open, write и close}{code:p3_write}{C}
#include <fcntl.h>
#include <stdio.h>
#include <unistd.h>
#include <sys/stat.h>
#include <sys/types.h>

int main() 
{
	struct stat statbuf;
	int fd1 = open("q.txt",O_RDWR);
	int fd2 = open("q.txt",O_RDWR);
	stat("q.txt", &statbuf);
	printf("inode: %lu;size: %ld; blksize: %ld\n", statbuf.st_ino, statbuf.st_size, statbuf.st_blksize);
	
	for(char c = 'a'; c <= 'z'; c++)
	{
		if (c%2)
		write(fd1, &c, 1);
		else
		write(fd2, &c, 1);
		stat("q.txt", &statbuf);
		printf("inode: %lu;size: %ld; blksize: %ld\n", statbuf.st_ino, statbuf.st_size, statbuf.st_blksize);
	}
	close(fd1);
	close(fd2);
	return 0;
}
\end{mylisting}

\subsection{Результат работы программы}

\img{p3_write}{Результат работы программы с open, write и close}{pic:p3_write}

\section{Анализ работы программы}

Системный вызов open создаёт два дескриптора открытого файла в режиме <<чтения и записи>> (fd1, fd2). Соответствующие дескрипторам структуры file ссылаются на один и тот же inode файла <<q.txt>>. Ссылки на эти структуры помещаются в таблицу открытых процессом файлов под индексами 3 и 4, так как дескрипторы 0, 1 и 2 заняты стандартными потоками.

В первой итерации цикла системный вызов write для дескриптора fd1 записывает символ <<a>> в открытый файл, что приводит к увеличению на 1 значения поля f\_pos в соответствующей структуре file. Размер файла становится равным 1 байту.

Так как дескрипторам fd1 и fd2 соответствуют независимые структуры file, то для fd2 значение поля f\_pos остаётся равным 0. Поэтому на второй итерации цикла вызов write для fd2 заменит символ <<a>> в файле на символ <<b>>. Размер файла остаётся равным 1 байту, а значение f\_pos для fd2 увеличится на 1.

Поочерёдная запись символов в файл продолжается до конца цикла.

Таким образом, в файл будет записана только каждая вторая буква алфавита, начиная с <<b>>, и итоговый размер файла составит 13 байт, так как вызов write для fd2 будет заменять символы, записанные вызовом write для fd1.

\section{Схема связи структур}

Схема связи структур представлена на рисунке~\ref{pic:scheme_2}.

\section{Вариант программы с fopen, fprintf и fclose}

\subsection{Код программы}

\begin{mylisting}{Код программы с fopen, fprintf и fclose}{code:p3_write}{C}
#include <fcntl.h>
#include <stdio.h>
#include <sys/stat.h>
#include <sys/types.h>

int main() 
{
	struct stat statbuf;
	FILE *fs1 = fopen("q.txt", "w");
	FILE *fs2 = fopen("q.txt", "w");
	stat("q.txt", &statbuf);
	printf("inode: %lu;size: %ld; blksize: %ld\n", statbuf.st_ino, statbuf.st_size, statbuf.st_blksize);
	
	for(char c = 'a'; c <= 'z'; c++)
	{
		if (c%2)
		fprintf(fs1, "%c", c);
		else
		fprintf(fs2, "%c", c);
	}
	stat("q.txt", &statbuf);
	printf("inode: %lu;size: %ld; blksize: %ld\n", statbuf.st_ino, statbuf.st_size, statbuf.st_blksize);
	fclose(fs2);
	fclose(fs1);
	stat("q.txt", &statbuf);
	printf("inode: %lu;size: %ld; blksize: %ld\n", statbuf.st_ino, statbuf.st_size, statbuf.st_blksize);
	return 0;
}
\end{mylisting}

\subsection{Результат работы программы}

\img{p3_fprintf}{Результат работы программы с fopen, fprintf и fclose}{pic:p3_fprintf}

\subsection{Анализ работы программы}

Вызов библиотечной функции fopen откроет файл <<q.txt>> в режиме <<только для записи>> и вернёт два указателя на инициализированные структуры FILE (fs1, fs2). Во время вызова создаются две незваисимые структуры file с дескрипторами 3 и 4, значения которых будут записаны в поля \_fileno соотвествующих структур \_IO\_FILE. Эти структуры file ссылаются на один и тот же inode файла <<q.txt>>.

Вызов библиотечной функции fprintf для fs1 записывает символ <<a>> в буфер, что приводит к смещению в соответствующей структуре указателя \_IO\_WRITE\_PTR на 1 байт.

Вызов библиотечной функции fprintf для fs2 записывает символ <<b>> в буфер, что приводит к смещению в соответствующей структуре указателя \_IO\_WRITE\_PTR на 1 байт.

В результате работы цикла в буфере fs1 будет находиться каждая вторая буква алфавита, начиная с <<a>>, а в fs2 --- каждая вторая буква алфавита, начиная с <<b>>.

Вызов fclose для fs2 приводит к записи из буфера в открытый файл. Значение поля f\_pos для дескриптора 4 и размера файла становится 13. После этого дескриптор закрывается.

Так как были созданы две структуры открытых файлов, ссылающихся на один и тот же inode, а в структуре, соответствующей fs1, значение поля f\_pos осталось равным 0, то при вызове fclose для fs1 буфер записывается в открытый файл с начала, что приводит к замене последовательности символов буфера fs2 на последовательность символов буфера fs1. Размер файла остаётся неизменным (13 байт).

Таким образом, итоговое содержимое файла зависит от порядка вызовов fclose. В данном случае в файл будет записан каждый второй символ алфавита, начиная с буквы <<a>>, а размер файла составит 13 байт.

\subsection{Схема связи структур}

\img{3.pdf}{Схема связи структур в третьей программе с fopen, fprintf и fclose}{pic:scheme_3}